\FigureLayoutDeclare{img/2015/national_paper_1_liberal/q6_fig1.png}{5.0cm}{8bp 31bp 17bp 24bp}

\chapter{2015年全国I卷（文）}

\section{单选题}

\begin{problem}
已知集合 \(A = \{x \mid x = 3n + 2, n \in \N\}\)，\(B = \{6, 8, 10, 12, 14\}\)，则集合 \(A \cap B\) 中的元素个数为（\quad）

\choices
    {\(5\)}
    {\(4\)}
    {\(3\)}
    {\(2\)}
\end{problem}

\begin{answer}
D
\end{answer}

\begin{solution}
在集合 \(B\) 中逐个检验 \(x=3n+2\).当 \(x=6\)，\(10\)，\(12\) 时，\(x-2\) 不是 \(3\) 的倍数；而 \(8=3\times2+2\)，\(14=3\times4+2\)，所以
\[
A\cap B=\{8,14\}
\]
因此 \(A\cap B\) 中有 \(2\) 个元素，故选 D.
\end{solution}

\begin{problem}
已知点 \(A(0,1)\)，\(B(3,2)\)，向量 \(\overrightarrow{AC} = (-4, -3)\)，则向量 \(\overrightarrow{BC} =\)（\quad）

\choices
    {\((-7, -4)\)}
    {\((7, 4)\)}
    {\((-1, 4)\)}
    {\((1, 4)\)}
\end{problem}

\begin{answer}
A
\end{answer}

\begin{solution}
由 \(\overrightarrow{AC}=(-4,-3)\)，得
\[
C=A+\overrightarrow{AC}=(0,1)+(-4,-3)=(-4,-2)
\]
因此
\[
\overrightarrow{BC}=C-B=(-4,-2)-(3,2)=(-7,-4)
\]
故选 A.
\end{solution}

\begin{problem}
已知复数 \(z\) 满足 \((z - 1)\i = 1 + \i\)，则 \(z =\)（\quad）

\choices
    {\(-2 - \i\)}
    {\(-2 + \i\)}
    {\(2 - \i\)}
    {\(2 + \i\)}
\end{problem}

\begin{answer}
C
\end{answer}

\begin{solution}
由 \((z-1)\i=1+\i\)，两边同除以 \(\i\)，得
\[
z-1=\frac{1+\i}{\i}=(1+\i)(-\i)=1-\i
\]
所以 \(z=2-\i\)，故选 C.
\end{solution}

\begin{problem}
如果 \(3\) 个正整数可作为一个直角三角形三条边的边长，则称这 \(3\) 个数为一组勾股数，从 \(1\)，\(2\)，\(3\)，\(4\)，\(5\) 中任取 \(3\) 个不同的数，则这 \(3\) 个数构成一组勾股数的概率为（\quad）

\choices
    {\(\frac{3}{10}\)}
    {\(\frac{1}{5}\)}
    {\(\frac{1}{10}\)}
    {\(\frac{1}{20}\)}
\end{problem}

\begin{answer}
C
\end{answer}

\begin{solution}
从 \(1\)，\(2\)，\(3\)，\(4\)，\(5\) 中任取 \(3\) 个不同的数，共有
\[
\mathrm C_5^3=10
\]
种取法.三数构成直角三角形三边时，最大数的平方应等于另外两个数的平方和；在这些取法中只有 \((3,4,5)\) 满足 \(3^{2}+4^{2}=5^{2}\).因此所求概率为
\[
\frac{1}{10}
\]
故选 C.
\end{solution}

\begin{problem}
已知椭圆 \(E\) 的中心为坐标原点，离心率为 \(\frac{1}{2}\)，\(E\) 的右焦点与抛物线 \(C: y^2 = 8x\) 的焦点重合，\(A\)，\(B\) 是 \(C\) 的准线与 \(E\) 的两个交点，则 \(|AB| =\)（\quad）

\choices
    {\(3\)}
    {\(6\)}
    {\(9\)}
    {\(12\)}
\end{problem}

\begin{answer}
B
\end{answer}

\begin{solution}
抛物线 \(y^{2}=8x\) 可写成 \(y^{2}=4px\)，所以 \(p=2\)，其焦点为 \((2,0)\)，准线为 \(x=-2\).

椭圆的右焦点为 \((2,0)\)，故 \(c=2\).由离心率 \(e=c/a=1/2\)，得 \(a=4\)，从而
\[
b^{2}=a^{2}-c^{2}=16-4=12
\]
椭圆方程为 \(\dfrac{x^{2}}{16}+\dfrac{y^{2}}{12}=1\).令 \(x=-2\)，得
\(\frac{4}{16}+\frac{y^{2}}{12}=1\)，\(\Longrightarrow\)，\(y^{2}=9\).
故 \(A\)，\(B\) 的纵坐标分别为 \(3,-3\)，所以 \(|AB|=6\)，故选 B.
\end{solution}

\begin{problem}
《九章算术》是我国古代内容极为丰富的数学名著，书中有如下问题：“今有委米依垣内角，下周\(8\text{~尺}\)，高\(5\text{~尺}\). 问：积及为米几何？”其意思为：“在屋内墙角处堆放米（如图，米堆为一个圆锥的四分之一），米堆底部的弧长为\(8\text{~尺}\)，米堆的高为\(5\text{~尺}\)，问米堆的体积和堆放的米各为多少？”已知\(1\text{~斛}\)米的体积约为\(1.62\text{~立方尺}\)，圆周率约为\(3\)，估算出堆放的米约有（\quad）

\bitmapfigure[max width=0.40\linewidth]{img/2015/national_paper_1_liberal/q6_fig1.png}

\choices
    {\(14\text{~斛}\)}
    {\(22\text{~斛}\)}
    {\(36\text{~斛}\)}
    {\(66\text{~斛}\)}
\end{problem}

\begin{answer}
B
\end{answer}

\begin{solution}
设圆锥底面半径为 \(r\).底面是四分之一圆，弧长为 \(8\)，所以
\[
\frac14\cdot2\pi r=8
\]
取 \(\pi\approx3\)，得 \(r=\dfrac{16}{3}\).米堆是圆锥的四分之一，其体积为
\[
V=\frac14\cdot\frac13\pi r^{2}\cdot5
=\frac{1}{12}\cdot3\cdot\left(\frac{16}{3}\right)^{2}\cdot5
=\frac{320}{9}\text{~立方尺}
\]
于是堆放的米约为
\[
\frac{V}{1.62}=\frac{320}{9\times1.62}\approx21.95\approx22\text{~斛}
\]
故选 B.
\end{solution}

\begin{problem}
已知 \( \{a_{n}\} \) 是公差为 \(1\) 的等差数列，\( S_{n} \) 为 \( \{a_{n}\} \) 的前 \(n\) 项和. 若 \( S_{8}=4S_{4} \)，则 \( a_{10}= \)（\quad）

\choices
    {\(\frac{17}{2}\)}
    {\(\frac{19}{2}\)}
    {\(10\)}
    {\(12\)}
\end{problem}

\begin{answer}
B
\end{answer}

\begin{solution}
设等差数列首项为 \(a_{1}\)，公差为 \(1\).由前 \(n\) 项和公式，得
\(S_{8}=8a_{1}+\frac{8\cdot7}{2}=8a_{1}+28\)，\(S_{4}=4a_{1}+\frac{4\cdot3}{2}=4a_{1}+6\).
由 \(S_{8}=4S_{4}\)，得
\[
8a_{1}+28=4(4a_{1}+6)
\]
所以 \(a_{1}=\dfrac12\).因此
\[
a_{10}=a_{1}+9\times1=\frac{19}{2}
\]
故选 B.
\end{solution}

\begin{problem}
函数 \( f(x) = \cos(\omega x + \varphi) \) 的部分图象如图所示，则 \( f(x) \) 的单调递减区间为（\quad）

\bitmapfigure[max width=0.45\linewidth]{img/2015/national_paper_1_liberal/q8_fig1.png}

\choices
    {\(\left(k\pi -\frac{1}{4},k\pi +\frac{3}{4}\right)\)，\(k\in \Z\)}
    {\(\left(2k\pi -\frac{1}{4},2k\pi +\frac{3}{4}\right)\)，\(k\in \Z\)}
    {\(\left(k - \frac{1}{4},k +\frac{3}{4}\right)\)，\(k\in \Z\)}
    {\(\left(2k -\frac{1}{4},2k +\frac{3}{4}\right)\)，\(k\in \Z\)}
\end{problem}

\begin{answer}
D
\end{answer}

\begin{solution}
由图象可读出函数在 \(x=\dfrac14\) 和 \(x=\dfrac54\) 处相邻过零，且两处过零方向相反，所以相邻过零点间距为半个周期，函数周期为 \(2\)，从而 \(\omega=\pi\).又图象在 \(x=\dfrac14\) 处下降，故可取
\[
f(x)=\cos\left(\pi x+\frac{\pi}{4}\right)
\]
余弦函数在区间 \((2k\pi,2k\pi+\pi)\) 上单调递减，所以
\[
2k\pi<\pi x+\frac{\pi}{4}<2k\pi+\pi
\]
即
\(2k-\frac14<x<2k+\frac34\)，\(k\in\Z\).
故选 D.
\end{solution}

\begin{problem}
执行如图的程序框图，如果输入的 \(t = 0.01\)，则输出的 \(n =\)（\quad）

\texfigure[max width=0.34\linewidth]{img_repaint/2015/national_paper_1_liberal/q9_fig1.tex}

\choices
    {\(5\)}
    {\(6\)}
    {\(7\)}
    {\(8\)}
\end{problem}

\begin{answer}
C
\end{answer}

\begin{solution}
初始时 \(S=1\)，\(n=0\)，\(m=\dfrac12\).每循环一次先作 \(S\leftarrow S-m\)，再作 \(m\leftarrow\dfrac m2\)，\(n\leftarrow n+1\).前六次循环后的 \(S\) 依次为
\(\frac12\)，\(\ \frac14\)，\(\ \frac18\)，\(\ \frac1{16}\)，\(\ \frac1{32}\)，\(\ \frac1{64}\)
均大于 \(t=0.01\).第七次循环后
\[
S=\frac1{128}=0.0078125<0.01
\]
此时停止并输出 \(n=7\)，故选 C.
\end{solution}

\begin{problem}
已知函数 \( f(x) = \begin{cases}
2^{x - 1} - 2, & x \leqslant 1,\\
-\log_2(x + 1), & x > 1,
\end{cases} \) 且 \( f(a) = -3 \)，则 \(f(6-a)=\)（\quad）

\choices
    {\(-\frac{7}{4}\)}
    {\(-\frac{5}{4}\)}
    {\(-\frac{3}{4}\)}
    {\(-\frac{1}{4}\)}
\end{problem}

\begin{answer}
A
\end{answer}

\begin{solution}
当 \(a\leqslant1\) 时，由 \(f(a)=-3\) 得
\[
2^{a-1}-2=-3
\]
即 \(2^{a-1}=-1\)，无实数解.因此 \(a>1\)，并且
\(-\log_{2}(a+1)=-3\)，\(\Longrightarrow\)，\(a+1=2^{3}=8\)
所以 \(a=7\).于是 \(6-a=-1\leqslant1\)，从第一段取值，得
\[
f(6-a)=f(-1)=2^{-2}-2=-\frac74
\]
故选 A.
\end{solution}

\begin{problem}
圆柱被一个平面截去一部分后与半球（半径为\(r\)）组成一个几何体. 该几何体三视图中的正视图和俯视图如图所示. 若该几何体的表面积为 \(16 + 20\pi\)，则 \(r =\)（\quad）

\bitmapfigure[max width=0.58\linewidth]{img/2015/national_paper_1_liberal/q11_fig1.png}

\choices
    {\(1\)}
    {\(2\)}
    {\(4\)}
    {\(8\)}
\end{problem}

\begin{answer}
B
\end{answer}

\begin{solution}
由三视图可知，该几何体由半球和半个圆柱拼接而成，半球、圆柱的半径均为 \(r\)，半个圆柱的高为 \(2r\).外露表面包括半球曲面、半个圆柱侧面、一个圆形底面和截面矩形，因此
\[
S=2\pi r^{2}+\pi r\cdot2r+\pi r^{2}+2r\cdot2r
=5\pi r^{2}+4r^{2}
\]
由题设
\[
(5\pi+4)r^{2}=16+20\pi=4(5\pi+4)
\]
由于 \(r>0\)，所以 \(r=2\)，故选 B.
\end{solution}

\begin{problem}
设函数 \( y = f(x) \) 的图象与 \( y = 2^{x + a} \) 的图象关于直线 \( y = -x \) 对称，且 \( f(-2) + f(-4) = 1 \)，则 \( a =\)（\quad）

\choices
    {\(-1\)}
    {\(1\)}
    {\(2\)}
    {\(4\)}
\end{problem}

\begin{answer}
C
\end{answer}

\begin{solution}
曲线 \(y=2^{x+a}\) 上的点 \((u,2^{u+a})\) 关于 \(y=-x\) 对称后变为 \((-2^{u+a},-u)\).因此对 \(x<0\)，对称曲线满足
\[
f(x)=a-\log_{2}(-x)
\]
所以
\(f(-2)=a-1\)，\(f(-4)=a-2\).
由 \(f(-2)+f(-4)=1\)，得 \(2a-3=1\)，所以 \(a=2\)，故选 C.
\end{solution}

\section{填空题}

\begin{problem}
数列\(\{a_{n}\}\) 中 \(a_{1}=2\)，\(a_{n+1}=2a_{n}\)，\(S_{n}\) 为\(\{a_{n}\}\) 的前 \(n\) 项和，若 \(S_{n}=126\)，则 \(n=\)\fillinblank{}.
\end{problem}

\begin{answer}
\(6\)
\end{answer}

\begin{solution}
由 \(a_{n+1}=2a_n\) 及 \(a_1=2\)，数列为首项 \(2\)、公比 \(2\) 的等比数列，所以
\[
a_n=2^n
\]
于是
\[
S_n=2+2^2+\cdots+2^n=2(2^n-1)
\]
由 \(S_n=126\)，得
\(2(2^n-1)=126\)，\(\Longrightarrow\)，\(2^n=64=2^6\)
所以 \(n=6\).
\end{solution}

\begin{problem}
已知函数 \( f(x) = ax^3 + x + 1 \) 的图象在点 \((1, f(1))\) 处的切线过点 \((2, 7)\)，则 \( a = \)\fillinblank{}.
\end{problem}

\begin{answer}
\(1\)
\end{answer}

\begin{solution}
函数的导数为
\[
f'(x)=3ax^2+1
\]
且 \(f(1)=a+2\).在 \(x=1\) 处的切线方程为
\[
y-(a+2)=(3a+1)(x-1)
\]
该直线过点 \((2,7)\)，代入得
\(7-(a+2)=3a+1\)，\(\Longrightarrow\)，\(5-a=3a+1\)
所以 \(a=1\).
\end{solution}

\begin{problem}
若 \(x\)，\(y\) 满足约束条件 \(\begin{cases}
x + y - 2 \leqslant 0,\\
x - 2y + 1 \leqslant 0,\\
2x - y + 2 \geqslant 0,
\end{cases}\) 则 \(z = 3x + y\) 的最大值为\fillinblank{}.
\end{problem}

\begin{answer}
\(4\)
\end{answer}

\begin{solution}
三条边界直线分别为
\(x+y=2\)，\(x-2y+1=0\)，\(2x-y+2=0\).
它们两两相交于
\((1,1)\)，\((0,2)\)，\((-1,0)\)
且这三个点均满足约束条件，因此可行域是以它们为顶点的三角形.在三个顶点处计算 \(z=3x+y\)，得
\(z(1,1)=4\)，\(z(0,2)=2\)，\(z(-1,0)=-3\).
所以 \(z\) 的最大值为 \(4\).
\end{solution}

\begin{problem}
已知 \(F\) 是双曲线 \(C: x^{2} - \frac{y^{2}}{8} = 1\) 的右焦点，\(P\) 是 \(C\) 左支上一点，\(A(0,6\sqrt{6})\)，当 \(\triangle APF\) 周长最小时，该三角形的面积为\fillinblank{}.
\end{problem}

\begin{answer}
\(12\sqrt{6}\)
\end{answer}

\begin{solution}
双曲线可写为
\[
\frac{x^2}{1^2}-\frac{y^2}{(2\sqrt2)^2}=1
\]
故 \(a=1\)，\(c=\sqrt{1+8}=3\).设左焦点为 \(F'(-3,0)\)，则 \(F=(3,0)\).对左支上的点 \(P\)，双曲线定义给出
\[
PF'-PF=2
\]
于是三角形 \(APF\) 的周长为
\[
AP+PF+AF=AP+PF'+AF-2
\]
由三角不等式，\(AP+PF'\geqslant AF'\)，等号成立时 \(A\)，\(P\)，\(F'\) 共线且 \(P\) 在线段 \(AF'\) 上.因此周长最小时，\(P\) 是直线 \(AF'\) 与双曲线左支的交点.

直线 \(AF'\) 的方程为
\[
y=2\sqrt6(x+3)
\]
与双曲线联立，得
\(x^2-3(x+3)^2=1\)，\(\Longrightarrow\)，\((x+7)(x+2)=0\).
线段 \(AF'\) 上的交点为 \(P=(-2,2\sqrt6)\).于是
\(\overrightarrow{AP}=(-2,-4\sqrt6)\)，\(\overrightarrow{AF}=(3,-6\sqrt6)\)
故三角形面积为
\[
S_{\triangle APF}
=\frac12\left|(-2)(-6\sqrt6)-(-4\sqrt6)\cdot3\right|
=12\sqrt6
\]
\end{solution}

\section{解答题}

\begin{problem}
已知 \(a\)，\(b\)，\(c\) 分别是 \(\triangle ABC\) 内角 \(A\)，\(B\)，\(C\) 的对边，\(\sin^{2} B = 2 \sin A \sin C\).

\begin{enumerate}
\item 若 \(a = b\)，求 \( \cos B \)；
\item 若 \(B = 90^{\circ}\)，且 \(a = \sqrt{2}\)，求 \(\triangle ABC\) 的面积.
\end{enumerate}
\end{problem}

\begin{solution}
由正弦定理，取外接圆半径为 \(R\)，有
\(\sin A=\frac{a}{2R}\)，\(\sin B=\frac{b}{2R}\)，\(\sin C=\frac{c}{2R}\).
代入题设，得
\[
\left(\frac{b}{2R}\right)^2
=2\cdot\frac{a}{2R}\cdot\frac{c}{2R}
\]
所以
\[
b^2=2ac
\]

\begin{enumerate}
\item 当 \(a=b\) 时，\(b^2=2bc\).由于边长为正，得 \(b=2c\)，即 \(c=\dfrac b2\).由余弦定理，
\[
\cos B=\frac{a^2+c^2-b^2}{2ac}
=\frac{b^2+\frac{b^2}{4}-b^2}{2b\cdot\frac b2}
=\frac14
\]

\item 当 \(B=90^{\circ}\) 时，勾股定理给出 \(b^2=a^2+c^2\).结合 \(b^2=2ac\)，得
\(a^2+c^2=2ac\)，\(\Longrightarrow\)，\((a-c)^2=0\).
所以 \(c=a=\sqrt2\).因此
\[
S_{\triangle ABC}=\frac12ac
=\frac12\cdot\sqrt2\cdot\sqrt2=1
\]
\end{enumerate}
\end{solution}

\begin{problem}
如图，四边形 \(ABCD\) 为菱形，\(G\) 为 \(AC\) 与 \(BD\) 的交点，\(BE \perp\) 平面 \(ABCD\).

\begin{enumerate}
\item 证明：平面 \(AEC \perp\) 平面 \(BED\)；
\item 若 \(\angle ABC = 120^{\circ}\)，\(AE \perp EC\)，三棱锥 \(E - ACD\) 的体积为 \(\frac{\sqrt{6}}{3}\)，求该三棱锥的侧面积.
\end{enumerate}

\bitmapfigure[max width=0.68\linewidth]{img/2015/national_paper_1_liberal/q18_fig1.png}
\end{problem}

\begin{solution}
\begin{enumerate}
\item 菱形的对角线互相垂直，所以 \(AC\perp BD\).又因为 \(BE\perp\) 平面 \(ABCD\)，且 \(AC\) 在该平面内，所以 \(AC\perp BE\).直线 \(BD\) 与 \(BE\) 相交于 \(B\)，并且都在平面 \(BED\) 内，故 \(AC\perp\) 平面 \(BED\).由于 \(AC\) 在平面 \(AEC\) 内，所以平面 \(AEC\perp\) 平面 \(BED\).

\item 设菱形边长为 \(s\)，设 \(BE=h\).由 \(\angle ABC=120^\circ\)，得
\(AC=\sqrt3s\)，\(BD=s\).
因为 \(BE\perp\) 平面 \(ABCD\)，所以
\(AE^2=AB^2+BE^2=s^2+h^2\)，\(CE^2=BC^2+BE^2=s^2+h^2\).
又 \(AE\perp EC\)，故在直角三角形 \(AEC\) 中
\[
AC^2=AE^2+CE^2=2(s^2+h^2)
\]
代入 \(AC^2=3s^2\)，得 \(h^2=\dfrac{s^2}{2}\).

三角形 \(ACD\) 是菱形的一半，所以
\[
S_{\triangle ACD}=\frac14AC\cdot BD=\frac{\sqrt3}{4}s^2
\]
三棱锥 \(E-ACD\) 的高为 \(h=\dfrac{s}{\sqrt2}\)，于是
\[
V=\frac13S_{\triangle ACD}h
=\frac13\cdot\frac{\sqrt3}{4}s^2\cdot\frac{s}{\sqrt2}
=\frac{\sqrt6}{24}s^3
\]
由 \(V=\dfrac{\sqrt6}{3}\)，得 \(s^3=8\)，所以 \(s=2\)，\(h=\sqrt2\).

此时
\(AE=CE=\sqrt{s^2+h^2}=\sqrt6\)，\(ED=\sqrt{BD^2+BE^2}=\sqrt6\).
因此
\[
S_{\triangle EAC}=\frac12AE\cdot CE=3
\]
在等腰三角形 \(EDA\) 和 \(ECD\) 中，腰长为 \(\sqrt6\)，底边长为 \(s=2\)，高均为
\[
\sqrt{(\sqrt6)^2-1^2}=\sqrt5
\]
所以
\[
S_{\triangle EDA}=S_{\triangle ECD}=\frac12\cdot2\cdot\sqrt5=\sqrt5
\]
三棱锥 \(E-ACD\) 的侧面积为
\[
3+\sqrt5+\sqrt5=3+2\sqrt5
\]
\end{enumerate}
\end{solution}

\begin{problem}
某公司为确定下一年度投入某产品的宣传费，需了解年宣传费 \(x\)（单位：\(\text{千元}\)）对年销售量 \(y\)（单位：\(\mathrm{t}\)）和年利润 \(z\)（单位：\(\text{千元}\)）的影响. 对近 \(8\) 年的年宣传费 \(x_i\) 和年销售量 \(y_i\)（\(i = 1\)，\(2\)，\(\cdots\)，\(8\)）数据作了初步处理，得到下面的散点图及一些统计量的值.

表中 \(w_i = \sqrt{x_i}\)，\(\overline{w} = \frac{1}{8}\sum_{i=1}^{8}w_i\).

\begin{center}
\renewcommand{\arraystretch}{1.2}
\begin{tabular}{|c|c|c|c|c|}
\hline
\(\overline{x}\) & \(\overline{y}\) & \(\overline{w}\) & \(\displaystyle\sum_{i=1}^{8}(x_i-\overline{x})^2\) & \(\displaystyle\sum_{i=1}^{8}(w_i-\overline{w})^2\) \\
\hline
\(46.6\) & \(563\) & \(6.8\) & \(289.8\) & \(1.6\) \\
\hline
\end{tabular}

\vspace{0.4em}

\begin{tabular}{|c|c|}
\hline
\(\displaystyle\sum_{i=1}^{8}(x_i-\overline{x})(y_i-\overline{y})\) & \(\displaystyle\sum_{i=1}^{8}(w_i-\overline{w})(y_i-\overline{y})\) \\
\hline
\(1.469\) & \(108.8\) \\
\hline
\end{tabular}
\end{center}

\begin{enumerate}
\item 根据散点图判断，\(y = a + bx\) 与 \(y = c + d\sqrt{x}\) 哪一个适宜作为年销售量 \(y\) 关于年宣传费 \(x\) 的回归方程类型？（给出判断即可，不必说明理由）
\item 根据（1）的判断结果及表中数据，建立 \(y\) 关于 \(x\) 的回归方程；
\item 已知这种产品的年利润 \(z\) 与 \(x\)，\(y\) 的关系为 \(z = 0.2y - x\). 根据（2）的结果回答下列问题：
\begin{enumerate}
\item 年宣传费 \(x = 49\) 时，年销售量及年利润的预报值是多少？
\item 年宣传费 \(x\) 为何值时，年利润的预报值最大？
\end{enumerate}
\end{enumerate}

附：对于一组数据 \((u_{1},v_{1})\)，\((u_{2},v_{2})\)，\(\dots\)，\((u_{n},v_{n})\)，其回归直线 \(v = \alpha + \beta u\) 的斜率和截距的最小二乘估计分别为 \(\hat{\beta} = \frac{\sum_{i=1}^{n}(u_i - \bar{u})(v_i - \bar{v})}{\sum_{i=1}^{n}(u_i - \bar{u})^2}\)，\(\hat{\alpha} = \bar{v} - \hat{\beta}\bar{u}\).

\bitmapfigure[max width=0.80\linewidth]{img/2015/national_paper_1_liberal/q19_fig1.png}
\end{problem}

\begin{solution}
\begin{enumerate}
\item 由散点图可判断，\(y=c+d\sqrt{x}\) 更适合作为年销售量 \(y\) 关于年宣传费 \(x\) 的回归方程类型.

\item 令 \(w=\sqrt{x}\)，建立 \(y\) 关于 \(w\) 的线性回归方程 \(\hat y=\hat c+\hat d w\).由题给数据，
\(\sum_{i=1}^{8}(w_i-\bar w)(y_i-\bar y)=108.8\)，\(\sum_{i=1}^{8}(w_i-\bar w)^2=1.6\).
因此
\[
\hat d=\frac{108.8}{1.6}=68
\]
\[
\hat c=\bar y-\hat d\,\bar w=563-68\times6.8=100.6
\]
因此 \(y\) 关于 \(x\) 的回归方程为
\[
\hat y=100.6+68\sqrt{x}
\]

\item
\begin{enumerate}
\item 当 \(x=49\) 时，
\[
\hat y=100.6+68\sqrt{49}=576.6
\]
所以年销售量的预报值为 \(576.6\mathrm{~t}\)，年利润的预报值为
\[
\hat z=0.2\hat y-x=0.2\times576.6-49=66.32\text{~千元}
\]

\item 令 \(u=\sqrt{x}\geqslant0\)，则年利润的预报值为
\[
\hat z(x)=0.2(100.6+68\sqrt{x})-x
=20.12+13.6u-u^2
=66.36-(u-6.8)^2
\]
当且仅当 \(u=6.8\) 时取得最大值，所以
\[
x=u^2=6.8^2=46.24
\]
即年宣传费为 \(46.24\text{~千元}\) 时，年利润的预报值最大.
\end{enumerate}
\end{enumerate}
\end{solution}

\begin{problem}
设函数 \(f(x)=\e^{2x}-a\ln x\).

\begin{enumerate}
\item 讨论 \(f(x)\) 的导函数 \(f'(x)\) 的零点的个数；
\item 证明：当 \(a>0\) 时，\(f(x)\geqslant2a+a\ln\frac{2}{a}\).
\end{enumerate}
\end{problem}

\begin{solution}
函数定义域为 \(x>0\)，且
\[
f'(x)=2\e^{2x}-\frac{a}{x}
=\frac{2x\e^{2x}-a}{x}
\]

\begin{enumerate}
\item 当 \(a\leqslant0\) 时，\(2\e^{2x}-\dfrac{a}{x}>0\)，所以 \(f'(x)\) 没有零点.当 \(a>0\) 时，令
\[
g(x)=2x\e^{2x}
\]
在 \(x>0\) 上，
\[
g'(x)=2\e^{2x}(1+2x)>0
\]
且 \(g(x)\) 从 \(0\) 增长到 \(+\infty\).因此方程 \(g(x)=a\) 有且只有一个正根，故 \(f'(x)\) 有且只有一个零点.

\item 设 \(x_0\) 是 \(f'(x)\) 的唯一零点，则
\[
2x_0\e^{2x_0}=a
\]
由 \(g(x)\) 的单调性，\(f'(x)<0\)（\(0<x<x_0\)），\(f'(x)>0\)（\(x>x_0\)），所以 \(f(x_0)\) 是 \(f\) 的最小值.又
\[
f(x_0)-2a-a\ln\frac{2}{a}
=\e^{2x_0}-2a-a\ln\frac{2x_0}{a}
=\e^{2x_0}-2a+2ax_0
=\e^{2x_0}(2x_0-1)^2\geqslant0
\]
因此
\[
f(x)\geqslant f(x_0)\geqslant2a+a\ln\frac{2}{a}
\]
\end{enumerate}
\end{solution}

\begin{problem}
已知过点 \(A(0,1)\) 且斜率为 \(k\) 的直线 \(l\) 与圆 \(C: (x - 2)^2 + (y - 3)^2 = 1\) 交于 \(M\)，\(N\) 两点.

\begin{enumerate}
\item 求 \(k\) 的取值范围；
\item 若 \(\overrightarrow{OM}\cdot\overrightarrow{ON}=12\)，其中 \(O\) 为坐标原点，求 \(|MN|\).
\end{enumerate}
\end{problem}

\begin{solution}
\begin{enumerate}
\item 直线 \(l\) 的方程为 \(y=kx+1\)，即 \(kx-y+1=0\).圆心为 \((2,3)\)，半径为 \(1\).要与圆交于两个不同的点，圆心到直线的距离应小于半径，因此
\[
\frac{|2k-2|}{\sqrt{k^2+1}}<1
\]
两边平方并整理，得
\[
3k^2-8k+3<0
\]
该二次式的两个根为
\(\frac{4-\sqrt7}{3}\)，\(\frac{4+\sqrt7}{3}\)
所以
\[
\frac{4-\sqrt7}{3}<k<\frac{4+\sqrt7}{3}
\]

\item 令 \(M=(x_1,kx_1+1)\)，\(N=(x_2,kx_2+1)\).将 \(y=kx+1\) 代入圆方程，得
\[
(1+k^2)x^2-4(1+k)x+7=0
\]
故
\(x_1+x_2=\frac{4(1+k)}{1+k^2}\)，\(x_1x_2=\frac{7}{1+k^2}\).
由向量数量积条件，
\[
\overrightarrow{OM}\cdot\overrightarrow{ON}
=x_1x_2+(kx_1+1)(kx_2+1)
\]
又
\[
(1+k^2)x_1x_2+k(x_1+x_2)+1
=8+\frac{4k(1+k)}{1+k^2}
\]
所以
\[
8+\frac{4k(1+k)}{1+k^2}=12
\]
所以 \(k=1\).此时直线为 \(y=x+1\)，代入圆方程得
\[
2x^2-8x+7=0
\]
于是
\[
|x_1-x_2|=\frac{\sqrt{(-8)^2-4\cdot2\cdot7}}{2}=\sqrt2
\]
两点在斜率为 \(1\) 的直线上，因此
\[
|MN|=|x_1-x_2|\sqrt{1+1^2}=\sqrt2\cdot\sqrt2=2
\]
\end{enumerate}
\end{solution}

\section{选考题：解答题}

\begin{problem}
如图，\(AB\) 是 \(\odot O\) 的直径，\(AC\) 是 \(\odot O\) 的切线，\(BC\) 交 \(\odot O\) 于点 \(E\).

\begin{enumerate}
\item 若 \(D\) 为 \(AC\) 的中点，证明：\(DE\) 是 \(\odot O\) 的切线；
\item 若 \(OA = \sqrt{3} CE\)，求 \(\angle ACB\) 的大小.
\end{enumerate}

\bitmapfigure[max width=0.36\linewidth]{img/2015/national_paper_1_liberal/q22_fig1.png}
\end{problem}

\begin{solution}
\begin{enumerate}
\item 先作相似变换令圆半径为 \(1\)，取
\(O=(0,0)\)，\(A=(-1,0)\)，\(B=(1,0)\)，\(C=(-1,t)\)
其中 \(t>0\).直线 \(BC\) 可参数表示为
\[
(x,y)=(1-2\lambda,t\lambda)
\]
与圆 \(x^2+y^2=1\) 的交点除 \(B\) 外对应
\((1-2\lambda)^2+t^2\lambda^2=1\)，\(\Longrightarrow\)，\(\lambda=\frac{4}{t^2+4}\).
所以
\(E=\left(\frac{t^2-4}{t^2+4},\frac{4t}{t^2+4}\right)\)，\(D=\left(-1,\frac t2\right)\).
直接计算得
\[
\overrightarrow{OE}\cdot\overrightarrow{OD}
=-\frac{t^2-4}{t^2+4}+\frac t2\cdot\frac{4t}{t^2+4}=1
\]
又 \(|OE|=1\)，所以
\[
\overrightarrow{OE}\cdot\overrightarrow{DE}
=\overrightarrow{OE}\cdot(\overrightarrow{OE}-\overrightarrow{OD})
=1-1=0
\]
即 \(OE\perp DE\)，故 \(DE\) 是 \(\odot O\) 在 \(E\) 处的切线.

\item 设 \(OA=r\)，\(AC=u\).由于半径垂直于切线，\(\triangle ABC\) 在 \(A\) 处为直角三角形，且 \(AB=2r\).由点 \(C\) 的切割线定理，
\[
CA^2=CE\cdot CB
\]
题设 \(OA=\sqrt3\,CE\)，所以 \(CE=\dfrac r{\sqrt3}\)，从而
\[
CB=\frac{CA^2}{CE}=\frac{\sqrt3u^2}{r}
\]
又由勾股定理，
\[
CB^2=u^2+4r^2
\]
令 \(v=\dfrac{u^2}{r^2}\)，得
\(3v^2=v+4\)，\(\Longrightarrow\)，\(3v^2-v-4=0\).
由于 \(v>0\)，所以 \(v=\dfrac43\)，即 \(u=\dfrac{2r}{\sqrt3}\).因此
\[
\tan\angle ACB=\frac{AB}{AC}=\frac{2r}{u}=\sqrt3
\]
\(\angle ACB\) 为直角三角形的锐角，故 \(\angle ACB=60^\circ\).
\end{enumerate}
\end{solution}

\begin{problem}
在直角坐标系 \(xOy\) 中，直线 \(C_1: x = -2\)，圆 \(C_2: (x - 1)^2 + (y - 2)^2 = 1\)，以坐标原点为极点，\(x\) 轴的正半轴为极轴建立极坐标系.

\begin{enumerate}
\item 求 \(C_{1}\)，\(C_{2}\) 的极坐标方程；
\item 若直线 \(C_3\) 的极坐标方程为 \(\theta = \frac{\pi}{4}\)（\(\rho \in \R\)），设 \(C_2\) 与 \(C_3\) 的交点为 \(M\)，\(N\)，求 \(\triangle C_2MN\) 的面积.
\end{enumerate}
\end{problem}

\begin{solution}
\begin{enumerate}
\item 由 \(x=\rho\cos\theta\)，\(y=\rho\sin\theta\)，直线 \(C_1\) 的极坐标方程为
\[
\rho\cos\theta=-2
\]
将 \(x=\rho\cos\theta\)、\(y=\rho\sin\theta\) 代入圆 \(C_2\)，得
\[
(\rho\cos\theta-1)^2+(\rho\sin\theta-2)^2=1
\]
即
\[
\rho^2-2\rho\cos\theta-4\rho\sin\theta+4=0
\]

\item 在 \(C_3\) 上有 \(\theta=\dfrac{\pi}{4}\)，故 \(x=y=\dfrac{\rho}{\sqrt2}\).代入 \(C_2\) 得
\[
\left(\frac{\rho}{\sqrt2}-1\right)^2+
\left(\frac{\rho}{\sqrt2}-2\right)^2=1
\]
整理为
\[
\rho^2-3\sqrt2\,\rho+4=0
\]
所以两个交点对应的 \(\rho\) 值为
\(\rho_1=\sqrt2\)，\(\rho_2=2\sqrt2\)
从而 \(|MN|=\rho_2-\rho_1=\sqrt2\).圆 \(C_2\) 的圆心为 \(Q=(1,2)\)，其到直线 \(y=x\) 的距离为
\[
\frac{|1-2|}{\sqrt2}=\frac1{\sqrt2}
\]
因此
\[
S_{\triangle QMN}
=\frac12\cdot\sqrt2\cdot\frac1{\sqrt2}
=\frac12
\]
\end{enumerate}
\end{solution}

\begin{problem}
已知函数 \( f(x) = |x + 1| - 2|x - a| \)，\( a > 0 \).

\begin{enumerate}
\item 当 \(a = 1\) 时，求不等式 \(f(x) > 1\) 的解集；
\item 若 \( f(x) \) 的图象与 \( x \) 轴围成的三角形面积大于 \(6\)，求 \( a \) 的取值范围.
\end{enumerate}
\end{problem}

\begin{solution}
\begin{enumerate}
\item 当 \(a=1\) 时，按 \(x<-1\)、\(-1\leqslant x<1\)、\(x\geqslant1\) 分段：
\[
f(x)=
\begin{cases}
x-3,&x<-1,\\
3x-1,&-1\leqslant x<1,\\
-x+3,&x\geqslant1.
\end{cases}
\]
第一段中 \(x-3>1\) 无解；第二段给出 \(x>\dfrac23\)；第三段给出 \(x<2\).合并得
\[
\frac23<x<2
\]

\item 对一般的 \(a>0\)，在断点 \(-1\) 和 \(a\) 处分段：
\[
f(x)=
\begin{cases}
x-(2a+1),&x<-1,\\
3x+1-2a,&-1\leqslant x<a,\\
-x+2a+1,&x\geqslant a.
\end{cases}
\]
函数与 \(x\) 轴的两个交点为
\(x_1=\frac{2a-1}{3}\)，\(x_2=2a+1\).
且 \(x_1<a<x_2\)，函数在两交点之间为正，在外部为负；最高点为 \((a,a+1)\).所以围成的三角形面积为
\[
S=\frac12(x_2-x_1)(a+1)
=\frac12\cdot\frac{4(a+1)}3\cdot(a+1)
=\frac{2(a+1)^2}{3}
\]
由 \(S>6\) 且 \(a>0\)，得
\((a+1)^2>9\)，\(\Longrightarrow\)，\(a>2\).
所以 \(a\in(2,+\infty)\).
\end{enumerate}
\end{solution}
